% SPDX-License-Identifier: Apache-2.0
% covers: EthShare
\datasheet{EthShare}{One Ethernet wire and two users: frames split by
EtherType on the way in, merged a frame at a time on the way out.}

\dsfacts{\code{lib/parts/src/ethshare.rs}}{\code{txhdl\_parts::ethshare::EthShare}}%
{\code{//docs:examples}}

\subsection*{Function}

The board has one Ethernet port and two things that want it: the
remote peripheral of issue 297, whose transport is a frame per
transaction, and the port a Zephyr driver talks to through
\code{EthSlots}. Neither can have the wire to itself.

Nothing new has to be invented to share it, because Ethernet already
carries several protocols on one wire and says which is which: bytes
twelve and thirteen of every frame are its EtherType. The remote
peripheral sends \code{0x88b5}, which IEEE leaves for local use, and
Zephyr's stack sends only the standard types, so the two cannot be
confused.

Its three inputs are the wire and the two users' frames to send, and
its three outputs the wire and the two users' frames received. The
first user is the one whose frames are of type \code{TYPE}; the
second takes every other frame.

\subsection*{Parameters}

\code{TYPE}, the first user's EtherType. The tree uses
\code{0x88b5}, the remote peripheral's, which
\code{remote::eth::ETHERTYPE} states; the tables are generated for
that.

\dstables{EthShare}

\subsection*{Behaviour}

\begin{itemize}
\item \textbf{Receiving holds a frame's first fourteen bytes before
either user sees any of it.} A frame's user is not known until its
fourteenth byte, and the thirteen before it belong to that user too.
So the unit holds them, reads the type from the last two, then hands
the held bytes on and passes the rest of the frame through to the
same user. Every frame received waits for its fourteenth byte.
\item Holding does not wait for either user. Handing on does: a user
that is not ready stops the frame, and the wire behind it.
\item \textbf{A frame that ends before its type is known goes to the
second user}, not to nobody. The receiver checks every frame and pads
a short one, so this should not arrive; but waiting for a fourteenth
byte that never comes would hold the wire for ever, which is worse
than delivering a frame nobody wants.
\item A frame of exactly fourteen bytes learns its user and its end
in the same cycle, and goes to the user its type names.
\item \textbf{Sending takes turns only when both users offer.} A frame
cannot be interleaved with another, so the two are merged a frame at
a time. When both offer, the one that did not send last goes first,
so neither can starve the other whichever sends more. A user with the
wire to itself is not slowed.
\end{itemize}

\subsection*{Verification}

\begin{itemize}
\item \code{ex\_ethshare} sends five frames in on the wire: two of the
first user's type and one of another, which is the split; one that
ends after ten bytes, before its type is known; and one of exactly
fourteen bytes. It checks that each user got its own frames whole and
in order. Both users also offer two frames each in the same cycle,
and the run checks that the four go out whole, taking turns.
\item The short frame is adversarial on purpose. Its last byte is
\code{0xb5}, the low half of the first user's type, and the frame
before it left \code{0x88} as the high half, so a unit that read the
type without checking that the frame had reached it would see
\code{0x88b5} and send the frame to the wrong user. The first version
of the run used bytes that did not form that type, and it passed with
the check removed.
\item The netlist is checked against that run under nvc and
Verilator: \code{//docs:ethshare\_sim\_ethshare\_tb\_test} and
\code{//docs:ethshare\_vsim\_test}.
\item On the board, both board tests of the Ethernet port in
\code{cpu/vreteno/tests/board.rs} go through it, one in each
direction.
\end{itemize}

\subsection*{Used by}

\code{Board}, as \code{share}, an \code{EthShare<0x88b5>} between the
wire and its two users: the remote peripheral's link first and the
Ethernet port's \code{FrameLen} and \code{FrameOut} second. The
example \code{ex\_ethshare}.

\subsection*{Limits}

It checks the type and nothing else. It does not look at addresses, so
both users see every frame of their type, whoever it was for.

It holds sixteen bytes of header memory and uses fourteen.

It runs on one clock, \code{DefaultClock}.
